\centerline{\bf \Huge Методы неравенств I}

\begin{flushright}
        \scriptsize{Лол}
\end{flushright}

\problem{1}
\textbf{Выпуклость.}
Докажите, что при $x, y, z \in[0,1]$ выполнено неравенство $x^2+y^2+z^2 \leqslant x y z+2$.

\problem{2}
\textbf{Неравенство о среднем.}
Положительные числа $a, b, c$ удовлетворяют соотношению $a b c=1$. Докажите, что $a^2+b^2+c^2 \geqslant a+b+c$.

\problem{3}
\textbf{Гомогенизация + Мюрхед.}
Для положительных $a, b, c$ верно, что $a+b+c=a b+b c+c a$, докажите, что

$$
a+b+c+1 \geqslant 4 a b c .
$$

\problem{4}
\textbf{Выравнивание знаменателей.}
Даны натуральное число $n>3$ и положительные числа $x_1, x_2, \ldots, x_n$, произведение которых равно 1 . Докажите неравенство

$$
\frac{1}{1+x_1+x_1 x_2}+\frac{1}{1+x_2+x_2 x_3}+\ldots+\frac{1}{1+x_n+x_n x_1}>1 .
$$

\problem{5}
\textbf{Угадайте сами!}
Для произвольных $a \geqslant 1$ и $b \geqslant 1$ докажите неравенство $2^{a b}-1 \geqslant\left(2^a-1\right)\left(2^b-1\right)$.

\problem{6}
\textbf{Подарок Дарине.}
Точка $O$ - центр описанной окружности треугольника $A B C$, а $E$ и $F$ - основания его высот, проведённых из вершин $B$ и $C$ соответственно. Найдите отношение площадей треугольников $B O F$ и $C O E$.